% Generated draft: requires mathematical and source review.
\documentclass[11pt]{article}
\usepackage[margin=25mm]{geometry}
\usepackage{fontspec}
\setmainfont{Latin Modern Roman}
\usepackage{amsmath,amssymb,mathtools,unicode-math}
\setmathfont{Latin Modern Math}
\usepackage{enumitem,needspace,xurl,hyperref}
\setlength{\parindent}{0pt}
\setlength{\parskip}{5pt}
\setlength{\emergencystretch}{3em}
\newcommand{\N}{\mathbb{N}}
\newcommand{\Z}{\mathbb{Z}}
\newcommand{\Q}{\mathbb{Q}}
\newcommand{\R}{\mathbb{R}}
\newcommand{\C}{\mathbb{C}}
\newcommand{\F}{\mathbb{F}}
\newcommand{\A}{\mathbb{A}}
\newcommand{\PP}{\mathbb{P}}
\newcommand{\E}{\mathbb{E}}
\newcommand{\eps}{\varepsilon}
\newcommand{\dd}{\,\mathrm{d}}
\begin{document}
\section*{UnsolvedMath ID: 30003446}
\textbf{Source problem code:} catalogue entry
\subsection{Definitions and mathematical statement}
Write \(\beta=\beta^{(2d)}\) for the supplied truncated moment data. The record does not specify the underlying space, the moment indexing convention, the permitted measures, or the meaning of a convergent extension sequence. Consequently, the following expresses the intended comparison conditionally; it is not a complete formulation of those missing definitions.

Once the intended moment setting is fixed, let \(\mathcal R(\beta)\) be its class of representing measures. Thus, if its underlying space is \(X\), its specified moment functions are \(\phi_\alpha\), and its truncated index set is \(I_{2d}\), a permitted measure \(\mu\) represents \(\beta\) when all required integrals exist and
\[
\int_X\phi_\alpha(x)\,\mathrm d\mu(x)=\beta_\alpha
\qquad(\alpha\in I_{2d}).
\]
These symbols are placeholders for conventions absent from the record, rather than a choice of a real, complex, or support-constrained moment problem.

Under a positive-measure interpretation, let \(\mathcal R_{\mathrm{fin}}(\beta)\) consist of the finitely atomic representing measures. For
\[
\mu=\sum_{j=1}^{m}w_j\delta_{x_j},
\qquad w_j>0,
\]
with distinct points \(x_j\), define \(a(\mu)=m\). Here \(\delta_x\) denotes unit mass at \(x\); the empty sum has zero atoms if the zero measure is permitted. When \(\mathcal R_{\mathrm{fin}}(\beta)\ne\varnothing\), define
\[
a_{\min}(\beta)=\min_{\nu\in\mathcal R_{\mathrm{fin}}(\beta)}a(\nu).
\]
The positivity convention and applicability of this finite-atomic comparison require confirmation from the source.

Let \(\mathcal C(\beta)\) denote the intended class of convergent extension sequences, and let \(\ell(E)\) denote the intended length of \(E\). If a shortest sequence exists, write
\[
\ell_{\min}(\beta)=\min_{F\in\mathcal C(\beta)}\ell(F).
\]
The supplied question asks whether a shortest convergent extension sequence necessarily produces a representing measure with the minimum possible number of atoms. On the universal reading, its conditional target is
\[
\ell(E)=\ell_{\min}(\beta)
\quad\Longrightarrow\quad
 a(\mu_E)=a_{\min}(\beta),
\]
for every admissible \(\beta\), every shortest \(E\in\mathcal C(\beta)\), and every measure \(\mu_E\) counted as produced by \(E\). This display presupposes that the relevant minima exist and that the produced measures are finitely atomic. The source does not define admissibility, extension steps, convergence, length, or the production of a measure. It also does not settle whether the intended quantification concerns every shortest sequence and every output, or the existence of a suitable shortest sequence and output. These unresolved qualifications prevent an unconditional precise statement.
\subsection{Short English statement}
The question compares two kinds of minimality: the length of an extension procedure for truncated moment data, and the number of atoms in a representing measure. Does minimizing the former necessarily minimize the latter? The supplied record leaves the moment setting, extension procedure, convergence criterion, and output quantifiers unspecified, so the displayed formulation is conditional and requires human mathematical review.
\subsection{Sources}
\begin{itemize}
\item Supplied local UnsolvedMath registry record: permanent ID \(30003446\), source problem code \texttt{catalogue entry}, title \emph{Shortest Moment Extensions and Minimal Atomic Measures}, in the Oberwolfach Reports Open Problems collection. No individual proposer is identified; the supplied proposed year is 2017.
\item Supplied report citation: \emph{Real Algebraic Geometry With a View Toward Moment Problems and Optimization} (2018), Oberwolfach Reports, pp.\ 771--862, DOI \texttt{10.4171/owr/2017/14}. \url{https://doi.org/10.4171/owr/2017/14}. The record describes Oberwolfach Reports as documenting mathematical workshops and their open-problem discussions. The exact problem location within the report is not supplied.
\item Supplied literature source, labelled ``Oberwolfach truncated-moment report'': \url{https://doi.org/10.4171/owr/2017/14}.
\item Supplied literature source, labelled ``Fialkow, truncated moment algorithms'': \url{https://jot.theta.ro/jot/archive/2008-060-002/2008-060-002-007.pdf}. A full bibliographic citation and a precise supporting passage are not supplied.
\item Literature assessment dated exactly \texttt{2026-08-21}: the supplied curation labels the problem open and states that bounds on lengths of convergent extension sequences and algorithms for special finite-variety cases are known, while the shortest-sequence versus fewest-atoms implication remains an explicit open question. This is supplied dated curation, not an independently verified assessment. Its background expressly requests independent verification of time-sensitive status and cited claims.
\item The supplied statement status is \texttt{corrected\_verified}; its accompanying explanation reports repairs to notation, typography, and sentence structure while preserving mathematical content. This records the supplied editorial assessment, not independent mathematical verification.
\end{itemize}
\end{document}
